\documentclass[11pt]{article}
\usepackage[utf8]{inputenc}
\usepackage{amsmath,amssymb,amsthm}
\usepackage[margin=1in]{geometry}
\usepackage{hyperref}
\usepackage{xcolor}
\usepackage{newunicodechar}
\newunicodechar{ℝ}{\ensuremath{\mathbb{R}}}
\newunicodechar{⟨}{\ensuremath{\langle}}
\newunicodechar{⟩}{\ensuremath{\rangle}}
\newunicodechar{·}{\ensuremath{\cdot}}
\newunicodechar{≠}{\ensuremath{\neq}}
\title{Marginal expectation and free-energy additivity for separable Gibbs measures}
\author{Tide retrospective, laplace seabed}
\date{2026-09-06}
\begin{document}
\sloppy
\maketitle

\section{Setting}
The laplace seabed's separable-potential module \texttt{TwoD/AddSeparable.lean} establishes the
factorisation facts for a product potential $L(x,y) = U(x)+V(y)$: the partition function factors,
$Z_{2D} = Z_U\cdot Z_V$; separable observables factor in expectation; and the mixed covariance
vanishes. A recent companion, \texttt{SeparableSameCoordCov}, added the same-coordinate covariance
inheritance at the second-cumulant level. This tide fills in the two remaining elementary facts of
the separable picture --- one at the mean level, one at the free-energy (CGF) level --- both of which
follow directly from the factorisation already in hand. It is a deliberate consolidation step rather
than a new frontier: the monomial cumulant ladder and the block-diagonal covariance structure being
complete, this rounds out the separable module's basic API.

\section{The thing formalised}
For a separable potential and both marginal partition functions nonzero, a single-coordinate
observable marginalises to its one-dimensional Gibbs expectation,
\[
  \langle f\circ\mathrm{fst}\rangle_{2D,t} = \langle f\rangle_{U,t},\qquad
  \langle g\circ\mathrm{snd}\rangle_{2D,t} = \langle g\rangle_{V,t},
\]
the mean-level companion of the covariance-level \texttt{SeparableSameCoordCov}. And the free energy
is additive,
\[
  \log Z_{2D}(t) = \log Z_U(t) + \log Z_V(t),
\]
the log-partition (cumulant-generating-function) reading of the partition-function factorisation. The
Lean names are \texttt{gibbsExpectation\_addSeparable\_fst} / \texttt{\_snd} and
\texttt{logPartition\_addSeparable}, in a new leaf file \texttt{TwoD/SeparableMarginal.lean}.

\section{Proof strategy}
The marginal expectation is the separable-observable factorisation
\texttt{gibbsExpectation\_separable\_addSeparable} run with the constant $1$ in the free coordinate:
$\langle f(z_1)\rangle_{2D} = \langle f\rangle_U\cdot\langle 1\rangle_V = \langle f\rangle_U$, since
the Gibbs expectation of the constant $1$ is $1$ (\texttt{gibbsExpectation\_const}). This is exactly
the collapse that was previously \emph{inlined} inside the cross-covariance-vanishing proof; the tide
lifts it to a named, reusable lemma. The free-energy additivity is one rewrite: the partition-function
factorisation followed by $\log(ab) = \log a + \log b$ (\texttt{Real.log\_mul}, whose two nonzero-base
side conditions are discharged by the positivity of each marginal partition function).

\section{Roadblocks and resolutions}
None --- both results compiled on the first attempt, and the full library rebuilt clean. As with the
sibling covariance tide, the only process choice was to place the additions in a new leaf file
importing \texttt{AddSeparable} rather than editing the core hub, so that no downstream file
re-elaborates and no latent missing-header warnings cascade. The build is warning-clean and
axiom-clean (the three standard axioms only).

\section{What was Mathlib, what was new}
The only Mathlib input is \texttt{Real.log\_mul}. Everything else is the seabed's own separable-%
factorisation API. What is new is the exposure of the marginal expectation as a first-class lemma
(previously reachable only by re-deriving the inlined step) and the CGF-level statement of free-energy
additivity, completing the elementary separable-potential toolkit at the mean and free-energy levels
alongside the existing partition, expectation, and covariance facts. No GPT consult was needed.
\end{document}
